\section{The scalar conversion}\label{app:scalar}

For completeness, we give the counting step that turns a chain
inequality into a hereditary bound on segment deficits. This is Pantone's scalar theorem~\cite{pantone-framework}; no
integrality of $p$ or $t$ is required.

Call a finite sequence of nonnegative integers \emph{legal} if its
zero intervals have length at most $a$, except that an entirely zero
sequence may have length $a+1$; every factor $(0,1,d)$ has $d\ge2$;
and across an internal zero interval of length $a$, the neighboring
positive entries sum to at least $a+2$. These are the local run rules
for genuine light segments. A bound on all legal sequences satisfying
$\CP(p,t)$ therefore applies to those segments.

\begin{lemma}\label{lem:scalar}
Let $a\ge4$, $0\le p\le t$, $t>0$, and $t(a-2)\ge4p+6$.
Set
\[
 \lambda=at-2,\qquad B=2t(a+1)+4p.
\]
Every legal sequence satisfying $\CP(p,t)$ on its zero-ended factors
has length $r$ and sum $D$ satisfying $\lambda D-2tr+B\ge0$.
\end{lemma}

\begin{proof}
First consider a zero-ended sequence with $s\ge1$ positive clusters.
Let cluster $j$ have length $m_j$ and sum $D_j$. The rule concerning
$(0,1,d)$ gives $D_j\ge m_j+1$, so $D_j-2\ge m_j-1\ge0$.
Moreover $D_j=2$ holds exactly for a regular cluster. Put
\[
 M=\sum_j(m_j-1),\quad V=\sum_j(D_j-2),\quad
 b=\#\{j:(m_j,D_j)\ne(1,2)\}.
\]
Write $u_0,\ldots,u_s$ for the zero-gap lengths, including the two end
gaps, and define
\[
 g_0=\#\{j:1\le j<s,\ u_j=a\},\quad
 H=\sum_{j=1}^{s-1}\max(0,a-1-u_j),\quad
 L=(a-u_0)+(a-u_s).
\]
Counting entries and deficit gives
\begin{equation}\label{eq:cluster-count}
 r=as+a+1+M-H+g_0-L,\qquad D=2s+V.
\end{equation}

Let $h_1$ count the gaps of length $a-2$ between originally adjacent
regular clusters, and let $h_2$ count the shorter such gaps. Cutting
at the $b$ other clusters and at these $h_2$ gaps leaves at most
$b+1+h_2$ paths of chains. The sum of their numbers of links is $h_1$.
Thus $\CP(p,t)$ gives
\[
 s-b\le p(b+1+h_2)+\frac t2h_1.
\]
Since $H\ge h_1+2h_2$ and $p\le t$, it follows that
\begin{equation}\label{eq:J}
 J:=4p+4(p+1)b+2tH-4s\ge0.
\end{equation}

Set $c=\lambda-2t$. The numerical hypothesis says $c\ge4(p+1)$.
Give cluster $j$ the reserve
\[
 \Phi_j=\lambda(D_j-2)-2t(m_j-1)
       -4(p+1)\ind_{(m_j,D_j)\ne(1,2)}.
\]
A regular cluster has reserve zero. Every other cluster has
$D_j-2\ge1$, and therefore
$\Phi_j\ge c(D_j-2)-4(p+1)\ge0$.

Suppose a zero gap of length $a$ separates clusters $i,j$.
The actual neighboring positive entries sum to at least $a+2$;
each other entry in those clusters contributes at least one. Hence
\[
 (D_i-2)+(D_j-2)\ge a-2+(m_i-1)+(m_j-1).
\]
It follows that
\[
 \Phi_i+\Phi_j\ge\lambda(a-2)-8(p+1)
 \ge\lambda(a-2)-2c
 =\lambda(a-4)+4t\ge4t.
\]

Form a graph on the positive clusters, with an edge across each gap
of length $a$. Its $g_0$ edges form a subgraph of a path. In particular,
every degree is at most two. Substituting \eqref{eq:cluster-count}
and \eqref{eq:J} gives the identity
\[
 \lambda D-2tr+B
 =J+2tL
 +\frac12\sum_{\{i,j\}\in E}(\Phi_i+\Phi_j-4t)
 +\frac12\sum_j(2-\deg(j))\Phi_j.
\]
Every term on the right is nonnegative.

An entirely zero sequence has length at most $a+1$, so its left side
is at least $4p$. An entirely positive sequence satisfies
$\lambda D-2tr\ge0$. In a mixed sequence, remove the initial and
final positive clusters. Each removed entry contributes at least
$\lambda-2t=c\ge0$. The remainder is either one zero interval or a
zero-ended sequence of the kind already considered. The empty sequence
also satisfies the claimed inequality.
\end{proof}
